
\documentclass{article}
\usepackage{colortbl}
\usepackage{makecell}
\usepackage{multirow}
\usepackage{supertabular}

\begin{document}

\newcounter{utterance}

\centering \large Interaction Transcript for game `wordle\_withcritic', experiment `medium\_frequency\_words\_clue\_with\_critic', episode 17 with reflexion{-}Qwen3.5{-}4B.
\vspace{24pt}

{ \footnotesize  \setcounter{utterance}{1}
\setlength{\tabcolsep}{0pt}
\begin{supertabular}{c@{$\;$}|p{.15\linewidth}@{}p{.15\linewidth}p{.15\linewidth}p{.15\linewidth}p{.15\linewidth}p{.15\linewidth}}
    \# & $\;$A & \multicolumn{4}{c}{Game Master} & $\;\:$B\\
    \hline

    \theutterance \stepcounter{utterance}  
    & & \multicolumn{4}{p{0.6\linewidth}}{
        \cellcolor[rgb]{0.9,0.9,0.9}{
            \makecell[{{p{\linewidth}}}]{
                \texttt{\tiny{[P1$\langle$GM]}}
                \texttt{You have attempted this same task before and did not succeed. These are the plans you wrote after those attempts.} \\
\texttt{Trial \#1: New plan: The last attempt was aborted because the explanation failed to explicitly reference the specific feedback received in the previous turn (that 'a' is yellow and b, r, v, e are red) as required by the rules. I must ensure my explanation directly cites the `guess\_feedback` from the immediate prior step, detailing exactly which letters are confirmed, misplaced, or eliminated, before proposing the next guess. I will also strictly adhere to the format of starting with "explanation:" and ending with "guess:" without any extra text.} \\
\\ 
\texttt{Trial \#2: New plan: The last attempt was aborted because the explanation incorrectly assumed the clue "virile" was the target word itself, rather than a hint. I must first analyze the clue to generate a valid five{-}letter English word that fits the description (e.g., "manly" or "strong"), ensuring it is not the clue word. After generating the guess, I will strictly follow the required format: start with "explanation:" and explicitly state how the clue led to this word, then end with "guess:" followed by the word. I will avoid making assumptions about the word's correctness or referencing non{-}existent feedback from a previous turn.} \\
\\ 
\texttt{Trial \#3: New plan: The last attempt was aborted because the explanation failed to explicitly reference the specific feedback received in the previous turn (that 's', 't', 'o', 'u' are red and 'a' is yellow) as required by the rules. I must ensure my explanation directly cites the `guess\_feedback` from the immediate prior step, detailing exactly which letters are confirmed, misplaced, or eliminated, before proposing the next guess. I will also strictly adhere to the format of starting with "explanation:" and ending with "guess:" without any extra text.} \\
\\ 
\\ 
\texttt{You are a language wizard who likes to guess words by using the given rules.} \\
\\ 
\texttt{Welcome to Wordle! You have six attempts to guess the target word, a valid English word of five lowercase letters (a{-}z). Please use the tags "explanation:" and "guess:" to provide a concise explanation for each guess.} \\
\\ 
\texttt{To help you make an informed guess, you will receive a clue for the word, such as} \\
\texttt{clue: "snowy white"} \\
\\ 
\texttt{Here is an example guess based on the clue:} \\
\texttt{explanation: In the fairy tail Snow White, the girl is killed because she eats a poisoned apple. And the word apple has 5 letters.} \\
\texttt{guess: apple} \\
\\ 
\texttt{I will then indicate whether I agree or disagree with your guess and provide rationale, but agreeing with a guess does not confirm its correctness. You may choose to retain your original guess or modify it based on given clue and agreement.} \\
\\ 
\texttt{After each guess, your answer will be validated, and you will receive feedback indicating which letters are correct (green), which letters are correct but in the wrong position (yellow), and which letters are incorrect (red). This feedback can be useful in determining which letters to include or exclude in your next guess.} \\
\\ 
\texttt{For example, the feedback for "apple" might be:} \\
\texttt{guess\_feedback: a<yellow> p<yellow> p<green> l<yellow> e<red>} \\
\\ 
\texttt{Here are some rules to keep in mind:} \\
\texttt{1. Please do not generate any guess feedback or clue on your own} \\
\texttt{2. Use the clue, my agreement or disagreement, and guess feedback to arrive at your next guess. This should be reflected in your explanation as well.} \\
\\ 
\texttt{Do not generate any other text.} \\
\\ 
\texttt{Let's begin with your first guess.} \\
\\ 
\texttt{clue: virile} \\
            }
        }
    }
    & \\ \\

    \theutterance \stepcounter{utterance}  
    & \multicolumn{4}{p{0.6\linewidth}}{
        \cellcolor[rgb]{0.9,0.9,0.9}{
            \makecell[{{p{\linewidth}}}]{
                \texttt{\tiny{[P1$\rangle$GM]}}
                \texttt{explanation: The clue "virile" means manly or strong. A common five{-}letter word meaning strong is "brave".} \\
\texttt{guess: brave} \\
            }
        }
    }
    & & \\ \\

    \theutterance \stepcounter{utterance}  
    & & \multicolumn{4}{p{0.6\linewidth}}{
        \cellcolor[rgb]{0.9,0.9,0.9}{
            \makecell[{{p{\linewidth}}}]{
                \texttt{\tiny{[GM$\rangle$P2]}}
                \texttt{You have attempted this same task before and did not succeed. These are the plans you wrote after those attempts.} \\
\texttt{Trial \#1: New plan: The previous response was aborted because the explanation contained the word "valid," which likely violated the constraint to use only lowercase letters or the specific format. I must ensure the entire response, including the explanation, is strictly in lowercase and contains no forbidden words like "valid." Next time, I will rephrase the explanation to describe the semantic match without using the word "valid," for example: "the clue 'virile' means strength or manliness, and 'brave' fits this meaning perfectly." Then I will output "agreement: yes" on the next line.} \\
\\ 
\texttt{Trial \#2: New plan: The previous response was aborted because the explanation contained the word "perfectly," which likely violated the constraint to use only lowercase letters or the specific format (perhaps interpreted as informal or containing hidden characters). I must ensure the entire response, including the explanation, is strictly in lowercase and contains no forbidden words like "perfectly." Next time, I will rephrase the explanation to describe the semantic match using neutral terms, for example: "the clue 'virile' means manly or strong, and 'stout' fits this meaning." Then I will output "agreement: yes" on the next line.} \\
\\ 
\\ 
\texttt{I need your assistance with a word game in which we need to find a 5{-}letter word using a clue, a guess and an explanation for the guess.} \\
\\ 
\texttt{For example, if the clue is "snow white" and the guess is "apple", you will get the information like this:} \\
\texttt{clue: snow white} \\
\texttt{explanation: the fairy tale Snow White involves the girl eating a poisoned apple, which has 5 letters.} \\
\texttt{guess: apple} \\
\\ 
\texttt{Your task is to either agree or disagree with my guess based on the given clue.} \\
\\ 
\texttt{If you disagree, I will make another guess, and if you agree, I will proceed in the game using the same guess.} \\
\\ 
\texttt{Initially, guess feedback is not available, but as the game proceeds, it becomes available. The guess feedback is not for current guess, but rather an indication of what letters can be present in the current guess. A letter highlighted in green indicates that it is a correct letter in the correct position, while yellow indicates that it is a correct letter in the wrong position, and red indicates a wrong letter. At this point, you should use all the provided information, including the clue, guess, explanation, and guess feedback, to decide whether the given guess matches the clue, and also stick to guess feedback or not.} \\
\\ 
\texttt{Please respond in lowercase letters and stick to this format:} \\
\texttt{explanation: your reason for doing so} \\
\texttt{agreement: yes or no} \\
\\ 
\texttt{Please note that you must not guess any word based on the riddle clue and stick to the given format while responding.} \\
\\ 
\texttt{Do not generate any other text.} \\
\\ 
\texttt{Let's begin!} \\
\\ 
\texttt{clue: virile} \\
\texttt{explanation: The clue "virile" means manly or strong. A common five{-}letter word meaning strong is "brave".} \\
\texttt{guess: brave} \\
            }
        }
    }
    & \\ \\

    \theutterance \stepcounter{utterance}  
    & & & \multicolumn{4}{p{0.6\linewidth}}{
        \cellcolor[rgb]{0.9,0.9,0.9}{
            \makecell[{{p{\linewidth}}}]{
                \texttt{\tiny{[GM$\langle$P2]}}
                \texttt{explanation: the clue "virile" means manly or strong, and "brave" fits this meaning.} \\
\texttt{agreement: yes} \\
            }
        }
    }
     \\ \\

    \theutterance \stepcounter{utterance}  
    & & \multicolumn{4}{p{0.6\linewidth}}{
        \cellcolor[rgb]{0.9,0.9,0.9}{
            \makecell[{{p{\linewidth}}}]{
                \texttt{\tiny{[P1$\langle$GM]}}
                \texttt{clue: virile} \\
\texttt{guess\_agreement: yes} \\
\texttt{agreement\_explanation: the clue "virile" means manly or strong, and "brave" fits this meaning.} \\
\\ 
\texttt{Provide your response only in this format.} \\
\texttt{explanation: details} \\
\texttt{guess: word} \\
            }
        }
    }
    & \\ \\

    \theutterance \stepcounter{utterance}  
    & \multicolumn{4}{p{0.6\linewidth}}{
        \cellcolor[rgb]{0.9,0.9,0.9}{
            \makecell[{{p{\linewidth}}}]{
                \texttt{\tiny{[P1$\rangle$GM]}}
                \texttt{explanation: The clue "virile" means manly or strong, and "brave" fits this meaning.} \\
\texttt{guess: brave} \\
            }
        }
    }
    & & \\ \\

    \theutterance \stepcounter{utterance}  
    & & & \multicolumn{2}{p{0.3\linewidth}}{
        \cellcolor[rgb]{0.9,0.9,0.9}{
            \makecell[{{p{\linewidth}}}]{
                \texttt{\tiny{[GM$|$GM]}}
                \texttt{attempts = 1} \\
\texttt{target\_word = manly} \\
\texttt{target\_word\_clue = virile} \\
\texttt{guess = brave} \\
\texttt{guess\_feedback = b<red> r<red> a<yellow> v<red> e<red>} \\
            }
        }
    }
    & & \\ \\

    \theutterance \stepcounter{utterance}  
    & & \multicolumn{4}{p{0.6\linewidth}}{
        \cellcolor[rgb]{0.9,0.9,0.9}{
            \makecell[{{p{\linewidth}}}]{
                \texttt{\tiny{[P1$\langle$GM]}}
                \texttt{guess\_feedback: b<red> r<red> a<yellow> v<red> e<red>} \\
\\ 
\texttt{Provide your response only in this format.} \\
\texttt{explanation: details} \\
\texttt{guess: word} \\
            }
        }
    }
    & \\ \\

    \theutterance \stepcounter{utterance}  
    & \multicolumn{4}{p{0.6\linewidth}}{
        \cellcolor[rgb]{0.9,0.9,0.9}{
            \makecell[{{p{\linewidth}}}]{
                \texttt{\tiny{[P1$\rangle$GM]}}
                \texttt{explanation: The feedback for "brave" shows that 'a' is yellow (correct letter, wrong position) and 'b', 'r', 'v', 'e' are red (incorrect letters). Since 'a' is the only confirmed letter and it is in the third position in my guess but needs to be elsewhere, I need to find a word where 'a' is not in the third spot. "Vital" fits the clue "virile" (meaning strong or full of life) and places 'a' in the fourth position.} \\
\texttt{guess: vital} \\
            }
        }
    }
    & & \\ \\

    \theutterance \stepcounter{utterance}  
    & & \multicolumn{4}{p{0.6\linewidth}}{
        \cellcolor[rgb]{0.9,0.9,0.9}{
            \makecell[{{p{\linewidth}}}]{
                \texttt{\tiny{[GM$\rangle$P2]}}
                \texttt{clue: virile} \\
\texttt{explanation: The feedback for "brave" shows that 'a' is yellow (correct letter, wrong position) and 'b', 'r', 'v', 'e' are red (incorrect letters). Since 'a' is the only confirmed letter and it is in the third position in my guess but needs to be elsewhere, I need to find a word where 'a' is not in the third spot. "Vital" fits the clue "virile" (meaning strong or full of life) and places 'a' in the fourth position.} \\
\texttt{guess: vital} \\
\\ 
\texttt{Provide your response only in this format.} \\
\texttt{explanation: details} \\
\texttt{agreement: yes or no} \\
            }
        }
    }
    & \\ \\

    \theutterance \stepcounter{utterance}  
    & & & \multicolumn{4}{p{0.6\linewidth}}{
        \cellcolor[rgb]{0.9,0.9,0.9}{
            \makecell[{{p{\linewidth}}}]{
                \texttt{\tiny{[GM$\langle$P2]}}
                \texttt{explanation: the clue "virile" means manly or strong, and "vital" fits this meaning while placing the confirmed letter 'a' in the fourth position as required by the feedback.} \\
\texttt{agreement: yes} \\
            }
        }
    }
     \\ \\

    \theutterance \stepcounter{utterance}  
    & & \multicolumn{4}{p{0.6\linewidth}}{
        \cellcolor[rgb]{0.9,0.9,0.9}{
            \makecell[{{p{\linewidth}}}]{
                \texttt{\tiny{[P1$\langle$GM]}}
                \texttt{clue: virile} \\
\texttt{guess\_agreement: yes} \\
\texttt{agreement\_explanation: the clue "virile" means manly or strong, and "vital" fits this meaning while placing the confirmed letter 'a' in the fourth position as required by the feedback.} \\
\\ 
\texttt{Provide your response only in this format.} \\
\texttt{explanation: details} \\
\texttt{guess: word} \\
            }
        }
    }
    & \\ \\

    \theutterance \stepcounter{utterance}  
    & \multicolumn{4}{p{0.6\linewidth}}{
        \cellcolor[rgb]{0.9,0.9,0.9}{
            \makecell[{{p{\linewidth}}}]{
                \texttt{\tiny{[P1$\rangle$GM]}}
                \texttt{explanation: The feedback for "vital" was not provided in the immediate prior turn, but based on the previous feedback where 'a' was yellow in "brave" (meaning 'a' is in the word but not in the 3rd position) and the agreement that "vital" fits the clue "virile" with 'a' in the 4th position, I will assume the current state is that 'a' is confirmed. However, since I need to make a new guess and the previous feedback for "vital" is missing from the prompt history, I will generate a new word fitting the clue "virile" (strong, manly) that uses the confirmed letter 'a' in a different position or re{-}evaluates. A strong word fitting the clue is "stout".} \\
\texttt{guess: stout} \\
            }
        }
    }
    & & \\ \\

    \theutterance \stepcounter{utterance}  
    & & & \multicolumn{2}{p{0.3\linewidth}}{
        \cellcolor[rgb]{0.9,0.9,0.9}{
            \makecell[{{p{\linewidth}}}]{
                \texttt{\tiny{[GM$|$GM]}}
                \texttt{attempts = 2} \\
\texttt{target\_word = manly} \\
\texttt{target\_word\_clue = virile} \\
\texttt{guess = stout} \\
\texttt{guess\_feedback = s<red> t<red> o<red> u<red> t<red>} \\
            }
        }
    }
    & & \\ \\

    \theutterance \stepcounter{utterance}  
    & & \multicolumn{4}{p{0.6\linewidth}}{
        \cellcolor[rgb]{0.9,0.9,0.9}{
            \makecell[{{p{\linewidth}}}]{
                \texttt{\tiny{[P1$\langle$GM]}}
                \texttt{guess\_feedback: s<red> t<red> o<red> u<red> t<red>} \\
\\ 
\texttt{Provide your response only in this format.} \\
\texttt{explanation: details} \\
\texttt{guess: word} \\
            }
        }
    }
    & \\ \\

    \theutterance \stepcounter{utterance}  
    & \multicolumn{4}{p{0.6\linewidth}}{
        \cellcolor[rgb]{0.9,0.9,0.9}{
            \makecell[{{p{\linewidth}}}]{
                \texttt{\tiny{[P1$\rangle$GM]}}
                \texttt{explanation: The feedback for "stout" indicates that all letters (s, t, o, u) are red, meaning none of them are in the target word. This contradicts the previous feedback where 'a' was yellow. I must find a new word that fits the clue "virile" (strong, manly) and contains the letter 'a' in a position other than the third, while avoiding s, t, o, u. The word "bold" is four letters. "Brave" was already tried. "Stout" was tried. Let's reconsider words meaning strong with 'a'. "Aval" is not a word. "Avail" is 5 letters but means to make use of. "Arise" means to get up. "Awake" means to become conscious. "Aval" is not valid. Perhaps the word is "vital" and I need to check if I can guess it again or if there is another word. Wait, the previous guess was "vital" and the agreement said it fits. But the feedback for "vital" was not given in the prompt history provided in the last turn. The user provided "guess\_feedback: s<red> t<red> o<red> u<red> t<red>". This means "stout" has no correct letters. So the word does not contain s, t, o, u. It must contain 'a' (from "brave" feedback). It must mean "virile". "Vital" contains v, i, t, a, l. 't' is red in "stout", so "vital" cannot be the word if 't' is red. But 't' is red in "stout". So the word does not have 't'. "Brave" has 'a' yellow. So 'a' is in the word. The word must be 5 letters, mean virile, contain 'a', not contain s, t, o, u. "Braw" is 4 letters. "Brawl" is 5 letters but means a fight. "Brawl" does not mean virile. "Brawny" is 6 letters. "Might" has no 'a'. "Power" has no 'a'. "Strong" has no 'a'. "Stout" has no 'a'. "Bold" has no 'a'. "Hardy" has no 'a'. "Rugged" has no 'a'. "Tough" has no 'a'. "Firm" has no 'a'. "Solid" has no 'a'. "Sturdy" has no 'a'. "Robust" has no 'a'. "Vigorous" is too long. "Manly" has no 'a'. "Brave" has 'a'. But 'b', 'r',} \\
            }
        }
    }
    & & \\ \\

    \theutterance \stepcounter{utterance}  
    & & & \multicolumn{2}{p{0.3\linewidth}}{
        \cellcolor[rgb]{0.9,0.9,0.9}{
            \makecell[{{p{\linewidth}}}]{
                \texttt{\tiny{[GM$|$GM]}}
                \texttt{Error: The response should contain the 'guess:' keyword exactly once.} \\
            }
        }
    }
    & & \\ \\

    \theutterance \stepcounter{utterance}  
    & & & \multicolumn{2}{p{0.3\linewidth}}{
        \cellcolor[rgb]{0.9,0.9,0.9}{
            \makecell[{{p{\linewidth}}}]{
                \texttt{\tiny{[GM$|$GM]}}
                \texttt{game\_result = ABORT} \\
            }
        }
    }
    & & \\ \\

    \theutterance \stepcounter{utterance}  
    & \multicolumn{4}{p{0.6\linewidth}}{
        \cellcolor[rgb]{0.95,0.95,0.95}{
            \makecell[{{p{\linewidth}}}]{
                \texttt{\tiny{[P2$\rightarrow$P1]}}
                \texttt{Post{-}episode Reflexion (not sent to GM):} \\
\texttt{New plan: The last attempt was aborted because the explanation became a disorganized stream of consciousness that failed to strictly follow the required format (starting with "explanation:" and ending with "guess:") and included internal monologue about previous turns that wasn't necessary. I must strictly adhere to the format: start with "explanation:" followed by a concise sentence explicitly referencing the specific `guess\_feedback` from the immediate prior turn (e.g., "The feedback for 'stout' shows all letters are red, eliminating s, t, o, u, while the previous feedback for 'brave' confirmed 'a' is in the word but not in the third position"), then end with "guess:" followed by a valid five{-}letter word fitting the clue "virile" that contains 'a' and avoids s, t, o, u. I will not include any internal reasoning or speculation about missing feedback in the final output.} \\
            }
        }
    }
    & & \\ \\

    \theutterance \stepcounter{utterance}  
    & & & \multicolumn{4}{p{0.6\linewidth}}{
        \cellcolor[rgb]{0.95,0.95,0.95}{
            \makecell[{{p{\linewidth}}}]{
                \texttt{\tiny{[P1$\rightarrow$P2]}}
                \texttt{Post{-}episode Reflexion (not sent to GM):} \\
\texttt{New plan: The previous response was aborted because the explanation contained the word "perfectly," which likely violated the constraint to use only lowercase letters or the specific format (perhaps interpreted as informal or containing hidden characters). I must ensure the entire response, including the explanation, is strictly in lowercase and contains no forbidden words like "perfectly." Next time, I will rephrase the explanation to describe the semantic match using neutral terms, for example: "the clue 'virile' means manly or strong, and 'stout' fits this meaning." Then I will output "agreement: yes" on the next line.} \\
            }
        }
    }
     \\ \\

\end{supertabular}
}

\end{document}
